\documentclass[12pt]{article}
\usepackage[utf8]{inputenc}
\usepackage[T1]{fontenc}
\usepackage{hyperref}
\title{Faliscan Language Resources for K-12 Curriculum Development}
\author{Compiled by Claude Sonnet 5.5 (claude-sonnet-5-5) for Andr\'as Kornai}
\date{}
\begin{document}
\maketitle

\section{Introduction}
Faliscan (ISO 639-3 xfa; Glottocode fali1291) is an extinct Italic language of the ager Faliscus in southern Etruria and northern Lazio, closely related to Latin. It is known from roughly 300 to 355 short inscriptions, mostly from the 7th to the 2nd centuries BC. The grammar can be described only partially. No native speakers exist, and I found no revitalization programme or school curriculum.

\section{Primary Educational Resources}
No primers or school materials were found.

\section{Linguistic and Reference Materials}
\begin{itemize}
\item Wikipedia overview: \url{https://en.wikipedia.org/wiki/Faliscan_language}
\item PRIN ``Italia Antica'' project (Florence), a digital project on ancient Italic languages that includes Faliscan: \url{https://www.prin-italia-antica.unifi.it/p149.html}
\item Rigobianco, overview of the Faliscan corpus and its digital (TEI-EpiDoc) treatment: \url{https://www.prin-italia-antica.unifi.it/upload/sub/immagini/Rigobianco-Digital-Italy.pdf}
\item Overview of Faliscan (Falisco), University of Venice repository: \url{https://iris.unive.it/retrieve/e4239ddd-ee09-7180-e053-3705fe0a3322/Falisco.pdf}
\item A doctoral thesis re-examining 44 Faliscan inscriptions: \url{https://tesidottorato.depositolegale.it/handle/20.500.14242/194697}
\item An Academia.edu item on Faliscan: \url{https://www.academia.edu/34993242}
\end{itemize}
Standard works mentioned in the search results are Bakkum, \emph{The Latin Dialect of the Ager Faliscus: 150 Years of Scholarship} (2009), and Giacomelli, \emph{La lingua falisca} (1963). I saw no URL for either.

\section{Religious and Community Resources}
None found.

\section{Digital and Audio Resources}
The digital corpus work above is the only digital resource found. No audio was found.

\section{Curriculum Development Context}
Faliscan is a corpus language. A classroom unit could use selected inscriptions alongside Latin, since the two are closely related. Reconstructed or spoken material would be a scholarly reconstruction, not a native tradition.
\end{document}